% Part: first-order-logic
% Chapter: sequent-calculus
% Section: quantifier-rules

\documentclass[../../../include/open-logic-section]{subfiles}

\begin{document}

\olfileid{fol}{seq}{qrl}

\olsection{పరిమాణీకరణ నియమాలు}

\subsection{$\lforall$కు నియమాలు}

\begin{defish}
\Axiom$ !A(t), \Gamma \fCenter \Delta$
\RightLabel{\LeftR{\lforall}}
\UnaryInf$ \lforall[x][!A(x)],\Gamma \fCenter \Delta$
\DisplayProof
\hfill
\Axiom$ \Gamma \fCenter \Delta, !A(a) $
\RightLabel{\RightR{\lforall}}
\UnaryInf$ \Gamma \fCenter \Delta, \lforall[x][!A(x)]$
\DisplayProof
\end{defish}

\LeftR{\lforall}లో $t$ ఒక సంవృత పదం (అంటే చరరాశులు లేని పదం).
\RightR{\lforall}లో $a$ ఒక \tetoken{స్థిరసంకేతం}{!!a{constant}}; అది \RightR{\lforall}
నియమంలోని కింది సీక్వెంట్‌లో ఎక్కడా కనిపించరాదు. \RightR{\forall}
నిగమనం యొక్క \emph{ఐగెన్ చరరాశి} అని $a$ను అంటాం.\footnote{పై
నియమంలో $a$ ఒక \tetoken{స్థిరసంకేతం}{!!a{constant}} అయినప్పటికీ చారిత్రక కారణాల వల్ల
“ఐగెన్ చరరాశి” అనే పేరునే ఉపయోగిస్తాం.}

\subsection{$\lexists$కు నియమాలు}

\begin{defish}
\Axiom$ !A(a), \Gamma \fCenter \Delta $
\RightLabel{\LeftR{\lexists}}
\UnaryInf$ \lexists[x][!A(x)], \Gamma \fCenter \Delta$
\DisplayProof
\hfill
\Axiom$ \Gamma \fCenter \Delta, !A(t) $
\RightLabel{\RightR{\lexists}}
\UnaryInf$ \Gamma \fCenter \Delta, \lexists[x][!A(x)]$
\DisplayProof
\end{defish}

ఇక్కడ కూడా $t$ ఒక సంవృత పదం; $a$ ఒక \tetoken{స్థిరసంకేతం}{!!a{constant}}, అది
\LeftR{\lexists} నియమంలోని కింది సీక్వెంట్‌లో కనిపించదు.
\LeftR{\lexists} నిగమనం యొక్క \emph{ఐగెన్ చరరాశి} అని $a$ను అంటాం.

\RightR{\lforall} లేదా \LeftR{\lexists} నిగమనంలోని కింది సీక్వెంట్‌లో
ఐగెన్ చరరాశి కనిపించకూడదనే షరతును \emph{ఐగెన్ చరరాశి షరతు} అంటాం.

\begin{explain}
\tetoken{ఒక సూత్రం}{!!a{formula}} అయిన $!A$లో \tetoken{చరం}{!!{variable}}~$x$ స్వేచ్ఛగా ఉంటే దానిని
$!A(x)$ అని రాస్తామనే
సంప్రదాయాన్ని గుర్తు చేసుకోండి. అదే సందర్భంలో $!A(t)$ అనేది
$\Subst{!A}{t}{x}$కు సంక్షిప్త రూపం. అందువల్ల \RightR{\lexists}
నియమాన్ని ఇలా కూడా రాయవచ్చు:
\begin{prooftree}
  \Axiom$\Gamma \fCenter \Delta, \Subst{!A}{t}{x}$
  \RightLabel{\RightR{\lexists}}
  \UnaryInf$\Gamma \fCenter \Delta, \lexists[x][!A]$
\end{prooftree}
$t$ ఇప్పటికే $!A$లో కనిపించవచ్చని గమనించండి; ఉదాహరణకు $!A$ అనేది
$\Atom{\Obj P}{t,x}$ కావచ్చు. కాబట్టి $\Gamma \Sequent \Delta,
\Atom{\Obj P}{t,t}$ నుంచి $\Gamma \Sequent \Delta,
\lexists[x][\Atom{\Obj P}{t,x}]$ను నిగమించడం \RightR{\lexists}ను
సరిగ్గా వర్తింపజేయడమే---పూర్వాధారంలోని $t$ ప్రత్యక్షతల్లో ఒకటి లేదా
అంతకంటే ఎక్కువను బద్ధ \tetoken{చరం}{!!{variable}}~$x$తో “మార్చవచ్చు”; అన్నిటినీ
మార్చాల్సిన అవసరం లేదు. అయితే \RightR{\lforall},
\LeftR{\lexists}లలోని ఐగెన్ చరరాశి షరతులు \tetoken{స్థిరసంకేతం}{!!{constant}}~$a$ అనేది
$!A$లో కనిపించకూడదని కోరుతాయి. అందువల్ల $\Gamma \Sequent \Delta,
\Atom{\Obj P}{a,a}$ నుంచి $\Gamma \Sequent \Delta,
\lforall[x][\Atom{\Obj P}{a,x}]$ను $\RightR{\lforall}$తో సరిగ్గా
నిగమించలేం.
\end{explain}

\begin{explain}
\RightR{\lexists}, \LeftR{\lforall}లలో $t$ అనే పదంపై ఏ పరిమితులూ
లేవు. మరోవైపు, \LeftR{\lexists}, \RightR{\lforall} నియమాల్లో ఐగెన్
చరరాశి షరతు ప్రకారం పై సీక్వెంట్‌లోని $!A(a)$ వెలుపల
\tetoken{స్థిరసంకేతం}{!!{constant}}~$a$ ఎక్కడా కనిపించరాదు. వ్యవస్థ నిర్దుష్టంగా ఉండేలా,
అంటే చెల్లుబాటు అయ్యే సీక్వెంట్లను మాత్రమే \tetoken{నిగమనం}{!!{derive}s} చేసేలా చేయడానికి
ఈ షరతు అవసరం. ఈ షరతు లేకపోతే కిందివి అనుమతించబడతాయి:
\begin{prooftree}
  \Axiom$!A(a) \fCenter !A(a)$
  \RightLabel{*\LeftR{\lexists}}
  \UnaryInf$\lexists[x][!A(x)] \fCenter !A(a)$
  \RightLabel{\RightR{\lforall}}
  \UnaryInf$\lexists[x][!A(x)] \fCenter \lforall[x][!A(x)]$
  \DisplayProof\bottomAlignProof
  \qquad
  \Axiom$!A(a) \fCenter !A(a)$
  \RightLabel{*\RightR{\lforall}}
  \UnaryInf$!A(a) \fCenter \lforall[x][!A(x)]$
  \RightLabel{\LeftR{\lexists}}
  \UnaryInf$\lexists[x][!A(x)] \fCenter \lforall[x][!A(x)]$
\end{prooftree}
కానీ $\lexists[x][!A(x)] \Sequent \lforall[x][!A(x)]$ చెల్లుబాటు కాదు.
\end{explain}

\end{document}
